\documentclass[10pt,a4paper]{amsart}
\usepackage[margin=19mm]{geometry}
\usepackage{amsmath,amssymb,mathtools}
\usepackage[scr=rsfs,cal=euler]{mathalfa}
\usepackage{xfrac}
\usepackage[unicode,hidelinks,bookmarks=false]{hyperref}
\newcommand{\ie}{i.e.\ }
\newcommand{\cf}{cf.\ }
\newcommand{\resp}{respectively\ }
\newcommand{\Subsection}{\subsection}
\providecommand{\nobreakdash}{\nobreak\hbox{-}}
\setlength{\emergencystretch}{2em}
\multlinegap0pt
\usepackage{bm}
\usepackage{bbm}
\usepackage{dsfont}
\usepackage{xspace}
\usepackage{cancel}
\usepackage{mathtools}
\usepackage{amsthm}
\makeatletter
\@ifundefined{lemma}{\newtheorem{lemma}{Lemma}}{}
\@ifundefined{definition}{\newtheorem{definition}[lemma]{Definition}}{}
\makeatother
\makeatletter
\newcommand{\Cov}{\operatorname{Cov}}
\newcommand{\N}{{\mathbb N}}
\expandafter\ifx\csname named\endcsname\relax
\newenvironment{named}[2]
	{\def\namedthmname{#1}
	\refstepcounter{namedthm}
	\namedthm[#2]\def\@currentlabel{#1}}
	{\endnamedthm}
\fi
\makeatother
\title[A pointwise ergodic theorem along return times of rapidly mi]{A pointwise ergodic theorem along return times of rapidly mixing systems}
\author{Sebasti\'an Donoso, Alejandro Maass, Vicente Saavedra-Araya}
\date{Source: arXiv:2502.17746v2 (CC BY 4.0)}
\begin{document}
\maketitle

\noindent\emph{Source and licence.} A pointwise ergodic theorem along return times of rapidly mixing systems. Version arXiv:2502.17746v2 (CC BY 4.0), distributed under the Creative Commons Attribution 4.0 International licence, as recorded in the arXiv licence metadata for this exact version and shown on the abstract page https://arxiv.org/abs/2502.17746v2. Full rights notice: licenses/arxiv-2502.17746-CC-BY-4.0.txt.\par\smallskip

\noindent\textit{Editorial scope.} Counted unit: the Lemma whose source label is \texttt{key} in the article (source0.tex lines 461--465, source label \texttt{key}), delivered together with the article's own complete original proof of it (source0.tex lines 467--509, one \texttt{proof} environment of 43 lines). No mathematics is written, repaired or re-derived: statement and proof are the article's own text line for line. Required context retained uncounted: P (lines 385--389). Every definition, display and label of those lines is retained unchanged. Omitted from this package and not redistributed: 1--384, 390--460, 466--466, 510--756. Nothing inside a retained excerpt is omitted or altered: every retained body is delivered exactly as the article's lines are written, so no omission receipt of any kind accompanies this package. Citations: the retained excerpts carry no citation command at all, so no bibliography entry of the article is transcribed and no reference is redistributed. Reference adaptation: none was needed and none was made; every reference inside the delivered unit resolves inside it, so no unresolved reference or citation is produced. Printed slips kept literal: no printed word, symbol, hypothesis or number of the counted statement or of its proof was altered. Layout and class: the article's own class is \texttt{amsart} with packages including ulem, xcolor, amssymb, amsfonts, hyperref, cleveref, textcomp, graphicx, bbm, enumerate, inputenc, fontenc, eucal, etoolbox; the wrapper is the lane's pinned \texttt{amsart} editorial wrapper at 10pt on A4, which restates the theorem-like environments the retained text needs and adds the article's own macro definitions that the retained text needs (\texttt{\textbackslash Cov}, \texttt{\textbackslash N}), with extra wrapper packages (bm, bbm, dsfont, xspace, cancel, mathtools). The delivered numbering is the unit's own. Assets: no retained span contains a figure environment, a graphics-inclusion command, a PDF-page inclusion, a table or a TikZ picture; no image file is copied into this package, the package declares no asset dependency and no image file is redistributed with it. Licence: version arXiv:2502.17746v2 of this article is distributed under the Creative Commons Attribution 4.0 International licence, as recorded in the arXiv licence metadata for this exact version and shown on the abstract page https://arxiv.org/abs/2502.17746v2. A full rights notice is delivered at licenses/arxiv-2502.17746-CC-BY-4.0.txt. Characters: every retained excerpt is pure ASCII, so the delivered engine, XeLaTeX with its default Latin Modern text font, reports no missing character.
 \begin{definition} Let $(\Omega,\mathcal{F},\mathbb{P})$ be a probability space. We say that the collection of measurable sets $(E_n)_{n\in\N}$ satisfies the property \eqref{P} if there exists a decreasing sequence $\rho:\N\to \mathbb{R}_{+}$  such that $\displaystyle \lim_{n\to\infty}\rho(n)=0$ and  for every finite sequence of integers $1 \leq n_1<n_2<\cdots<n_k$,
 \begin{equation}
     \left|\mathbb{P}(E_{n_1} \cap \cdots \cap E_{n_k})-\mathbb{P}(E_{n_1})\mathbb{P}(E_{n_2}\cap \cdots \cap E_{n_k})\right|\leq \rho(n_2-n_1)\mathbb{P}(E_{n_2}\cap \cdots \cap E_{n_k}).  \tag{P}\label{P}
 \end{equation}
 \end{definition}

\begin{lemma} Let $(\Omega,\mathcal{F},\mathbb{P})$ be a probability space, and let $\mathcal{E}=(E_n)_{n\in \N}$ be a sequence of measurable sets in $\mathcal{F}$ of decreasing measure converging to zero and satisfying \eqref{P}. Let $X_n(\omega):=\mathbbm{1}_{E_n}(\omega)$, $\sigma_n:=\mathbb{E}_{\mathbb{P}}(X_n)$ and $Y_n(\omega):=X_n(\omega)-\sigma_n$. Then, for every $1\leq n_1\leq n_2\leq n_3\leq n_4$ such that $d:=n_2-n_1=n_4-n_3$, we can estimate the $4$-fold covariance as
\begin{align*}
    \mathbb{E}_{\mathbb{P}}\left( Y_{n_1} Y_{n_2} Y_{n_3} Y_{n_4} \right)&\ll  \sigma_{n_1}\rho(d)\Big(\sigma_{n_1}+\rho(n_3-n_2)\Big).
\end{align*}\label{key}
\end{lemma}

\begin{proof}
  Let $1\leq n_1\leq n_2\leq n_3\leq n_4$ be integers such that $d:=n_4-n_3=n_2-n_1$. By expanding $\mathbb{E}_{\mathbb{P}}(Y_{n_1}Y_{n_2}Y_{n_3}Y_{n_4})$  and grouping elements together, we can write
\begin{equation}
\begin{split}
\mathbb{E}_{\mathbb{P}}\left( Y_{n_1} Y_{n_2} Y_{n_3} Y_{n_4} \right) &=\Cov(X_{n_1},X_{n_2}X_{n_3}X_{n_4})-\sigma_{n_4}\Cov(X_{n_1},X_{n_2}X_{n_3})\\
&\quad-\sigma_{n_3}\Cov(X_{n_1},X_{n_2}X_{n_4})+\sigma_{n_3}\sigma_{n_4}\Cov(X_{n_1},X_{n_2})\\
&\quad-\sigma_{n_2}\Cov(X_{n_1},X_{n_3}X_{n_4})+\sigma_{n_2}\sigma_{n_4}\Cov(X_{n_1},X_{n_3})\\
&\quad+\sigma_{n_2}\sigma_{n_3}\Cov(X_{n_1},X_{n_4}).
\end{split}  \label{split1} 
\end{equation}
Since $\Cov(X_{n_1},X_{n_2}X_{n_3}X_{n_4})=\mathbb{P}(E_{n_1}\cap E_{n_2}\cap E_{n_3}\cap E_{n_4})-\mathbb{P}(E_{n_1})\mathbb{P}(E_{n_2}\cap E_{n_3}\cap E_{n_4})$ and the collection $\mathcal{E}$ satisfies \eqref{P},
\begin{align*}
    |\Cov(X_{n_1},X_{n_2}X_{n_3}X_{n_4})|&\leq \rho(n_2-n_1)\mathbb{P}(E_{n_2}\cap E_{n_3}\cap E_{n_4})\\
      &=\rho(n_2-n_1)\Big(\Cov(X_{n_2},X_{n_3}X_{n_4})+\sigma_{n_2}\mathbb{P}(E_{n_3}\cap E_{n_4})\Big) \\
      &\leq \rho(n_2-n_1)\Big(\rho(n_3-n_2)\mathbb{P}(E_{n_3}\cap E_{n_4})+\sigma_{n_2}\mathbb{P}(E_{n_3}\cap E_{n_4})\Big) \\
      &\leq \rho(n_2-n_1)\Big(\rho(n_3-n_2)+\sigma_{n_2}\Big)\Big(\rho(n_4-n_3)+\sigma_{n_3}\Big)\sigma_{n_4}\\
      &\leq \rho(d)\Big(\rho(n_3-n_2)+\sigma_{n_1}\Big)\Big(\rho(d)+\sigma_{n_1}\Big)\sigma_{n_1},
\end{align*}
where in the last step we used that $(\sigma_n)_{n\in \N}$ is decreasing and $n_2-n_1=n_4-n_3=d$. Using that $\rho(n)$ and $\sigma_n$ are decreasing to 0, we conclude 
\[ |\Cov(X_{n_1},X_{n_2}X_{n_3}X_{n_4})|\ll \rho(d)\rho(n_3-n_2)\sigma_{n_1}+\rho(d)\sigma_{n_1}^2.\]Similarly, we obtain
\begin{equation}
\begin{split}
    &|\sigma_{n_4}\Cov(X_{n_1},X_{n_2}X_{n_3})|\leq \rho(d)\Big(\rho(n_3-n_2)+\sigma_{n_1}\Big)\sigma_{n_1}^2,\\
    &|\sigma_{n_3}\Cov(X_{n_1},X_{n_2}X_{n_4})|\leq \rho(d)\Big(\rho(n_4-n_2)+\sigma_{n_1}\Big)\sigma_{n_1}^2, \text{ and }\\
    &|\sigma_{n_2}\Cov(X_{n_1},X_{n_3}X_{n_4})|\leq \rho(n_3-n_1)\Big(\rho(d)+\sigma_{n_1}\Big)\sigma_{n_1}^2.\label{split3} 
\end{split}
\end{equation}
As $\rho$ is decreasing, each of the terms in \eqref{split3} can be bounded above by $$\rho(d)\rho(n_3-n_2)\sigma_{n_1}^2+\rho(d)\sigma_{n_1}^3.$$
Also, note that
\begin{equation}
\begin{split}
    &|\sigma_{n_3}\sigma_{n_4}\Cov(X_{n_1},X_{n_2})|\leq \rho(d)\sigma_{n_1}^3,\\
    &|\sigma_{n_2}\sigma_{n_4}\Cov(X_{n_1},X_{n_3})|\leq \rho(n_3-n_1)\sigma_{n_1}^3,\text{ and }\\
    &|\sigma_{n_2}\sigma_{n_3}\Cov(X_{n_1},X_{n_4})|\leq \rho(n_4-n_1)\sigma_{n_1}^3. \label{split4} 
\end{split}
\end{equation}
Hence, each term in \eqref{split4} can be bounded above by $\rho(d)\sigma_{n_1}^3.$
Replacing all the estimations in \eqref{split1}, it follows that
\begin{align*}
    \mathbb{E}_{\mathbb{P}}\left( Y_{n_1} Y_{n_2} Y_{n_3} Y_{n_4} \right)&\ll \rho(d)\rho(n_3-n_2)\sigma_{n_1}+\rho(d)\sigma_{n_1}^2\\
    &+\rho(d)\rho(n_3-n_2)\sigma_{n_1}^2+\rho(d)\sigma_{n_1}^3+\rho(d)\sigma_{n_1}^3\\
    &\ll \rho(d)\rho(n_3-n_2)\sigma_{n_1}+\rho(d)\sigma_{n_1}^2.
\end{align*}\end{proof}

\end{document}
